\section{Removing the analytical comparator cutoff}
\label{sec:zero-comparison}
The cutoff used in the auxiliary-prior proof need not remain an
assumption on the final neural-network oracle. A uniform zero-function
comparison deals with the complementary resource range.

\begin{lemma}[A width-independent zero small ball]
\label{lem:raw-zero-smallball}
For the Gaussian neuron-gated model with bounded hidden activation,
and for the normalized Dirichlet Gaussian readout, there are constants
$c,C>0$, independent of dimension, depth, and learner width, such that
\[
 \mu_N\{\|z\|_\infty\le\epsilon\}\ge c\epsilon^C,
 \qquad0<\epsilon\le1.
\]
For a general globally Lipschitz hidden activation in the unnormalized
model, the same assertion holds with constants allowed to depend on
fixed $d,H$. It also holds for the independent variance-gamma readout
with bounded hidden activations described below.
\end{lemma}
\begin{proof}
For the unnormalized Gaussian-gate readout, write the final gates as
$S_j=T P_j$, with $T\sim\operatorname{Gamma}(\kappa,1)$ independent
of the Dirichlet proportions, hidden features, and final Gaussian
factors. If $R_j=\sup_x|h_j(x)|$, then
\[
 \|z\|_\infty\le|G_0|+T Z,
 \qquad Z=\sum_jP_j|G_j|R_j.
\]
The expectation of $Z$ is bounded uniformly in width. For bounded
activations it is also independent of dimension and depth; for a
Lipschitz activation use the uniform first-feature moments already
proved. The event $Z\le2\E Z+1$ has probability at least $1/2$.
Independently require $T\le c\epsilon/(1+\E Z)$ and
$|G_0|\le\epsilon/2$. The gamma and Gaussian small-ball bounds prove
the assertion.

For the normalized readout, introduce a size-biased index $I$ with
$\Pr(I=j\mid P)=P_j$. When $N\ge2$, its mass has distribution
$P_I\sim\operatorname{Beta}(1+\kappa/N,\kappa-\kappa/N)$.
Direct integration near one gives
\[
 \Pr(P_I\ge1-\delta)\ge c_\kappa\delta^\kappa
 \qquad(0<\delta\le1/2),
\]
uniformly in $N$. Given this mass and index, the remaining normalized
proportions have their usual Dirichlet law. Their gamma-ratio moment
formula yields
\[
 \E\left[\sum_{j\ne I}\sqrt{P_j}|G_j|
                       \mid P_I\right]
       \le C_\kappa\sqrt{1-P_I}.
\]
Choose $\delta=c\epsilon^2$, require the inactive weighted sum to be
at most $\epsilon/[4(1+M_\phi)]$, and require both $|G_I|$ and
$|G_0|$ to be at most the same threshold. Markov's inequality gives
conditional probability at least $1/2$ for the weighted-sum event;
the two independent Gaussian intervals contribute at least
$c\epsilon^2$. The resulting raw output is at most $\epsilon$.
For $N=1$ only the two Gaussian intervals are needed. Introducing the
auxiliary index does not change the marginal prior, so this proves its
small-ball lower bound.

For variance-gamma readout increments, their total absolute Poisson
jump mass has a fixed-shape gamma distribution, independently of width.
The same is true after adding a fixed-intensity readout bias. Since
hidden features are bounded, conditioning this total variation to be
at most a constant times $\epsilon$ proves the claim.
\end{proof}

\begin{proposition}[A general cutoff-removal rule]
\label{prop:raw-cutoff-removal}
Suppose the raw-tail oracle has native moment exponent $r$, a uniform
zero small ball as above, and a comparison resource satisfying
\[
 K_n(g)\ge c\{1+\|g\|_\infty^{1/r}\}.
\]
Then, after enlarging the oracle constant, its bound
$R(g)+C n^{-1/(2r+1)}K_n(g)+C(n+1)^{-2}$ holds for every comparator,
without the condition $\|g\|_\infty\le M_n/2$.
\end{proposition}
\begin{proof}
Condition the native prior on the zero small ball at radius
$(n+1)^{-2}$. Its KL is $O(\log(n+1))$ and its uniform squared error
from zero is at most $(n+1)^{-4}$. Theorem~\ref{thm:raw-tail-oracle} with
comparison zero therefore gives
\[
 \E R(\bar z)\le\E Y^2+C n^{-1/(2r+1)}\log(n+1)+C(n+1)^{-2}
                  \le C_{K_Y}.
\]
If $\|g\|_\infty>M_n/2$, the displayed resource lower bound and
$M_n=c n^{r/(2r+1)}$ imply
$n^{-1/(2r+1)}K_n(g)\ge c'>0$.
Increasing its coefficient makes the desired oracle inequality follow
from the uniform risk bound and $R(g)\ge0$. The complementary case
is exactly the original raw-tail theorem.
\end{proof}

For bounded hidden activations, the unnormalized Gaussian comparison
resource dominates $1+\|g\|_\infty$, and hence also
$1+\|g\|_\infty^{2/3}$. The normalized readout resource dominates
the square of the total readout variation, including the bias, and
therefore dominates $\|g\|_\infty^2$. The variance-gamma resource is
linear in an upper bound for readout variation and dominates
$\|g\|_\infty$. Thus the cutoff-removal proposition applies to all
three principal raw oracles.

For the fixed-depth general Lipschitz comparison,
\[
 \|g\|_\infty\le C_{d,H}(m+1)^H(B+1)^{H+1},
 \qquad r=(3H+1)/2.
\]
Both exponents $H/r$ and $(H+1)/r$ are at most one. Hence its
Gaussian comparison resource also dominates
$1+\|g\|_\infty^{1/r}$, with a fixed $d,H$ constant. Its raw
sub-Gaussian-label oracle can likewise omit the comparator cutoff.
